\section{数据是智能边界：公域、领域与产品内生数据}

前两章讲能力如何被压缩出来，本章讨论能力边界由什么决定。王冠把“有多少人工就有多少智能”改写为“有多少数据就有多少智能”。这里的数据不只是文件或表格，而是对问题的理解、业务 know-how、任务过程和用户反馈。

\begin{figure}[H]
\centering
\includegraphics[width=0.96\textwidth]{data-three-stages.png}
\caption{数据三阶段：公域数据、领域数据、产品内生数据对应不同玩家。自制概念图，依据 00:37:11--01:05:36 对谈内容整理。}
\end{figure}

\begin{knowledgebox}{读图：每一阶段利好的玩家不同}
公域数据阶段利好基座模型公司，因为大家都在清洗和压缩互联网数据；领域数据阶段利好大厂和传统行业，因为它们拥有业务场景和私有数据；产品内生数据阶段才更像应用创业机会，因为新产品能在使用中产生之前不存在的数据。
\end{knowledgebox}

\subsection{智能的时空观：数据、算力和算法}

王冠用一个二维空间的类比解释智能边界。数据决定圆的边界，即模型最终能到达的能力范围；算力决定逼近边界的速度；算法决定逼近边界的路径和效率。这个说法不否认算法重要，而是把算法放回数据边界之内：没有相应数据，算法再强也很难凭空得到稳定能力。

\begin{importantbox}{数据决定智能边界}
如果一个任务需要行业 know-how、真实场景和开放性判断，那么“再搞一些数据”通常不是脏活，而是在定义问题本身。产品经理之所以重要，是因为他们能把业务理解转成可学习数据。
\end{importantbox}

\subsection{为什么前两个阶段不是典型应用创业机会}

本节把数据三阶段进一步映射到创业机会。不是每一种数据都适合创业公司获取，也不是每一个阶段都适合应用公司切入。

公域数据阶段比拼人才密度、算力、工程效率和组织消耗，更适合基座模型公司。领域数据阶段利好大厂和信息化程度高的行业，因为它们本来就有业务场景、客户渠道和私有数据。应用创业公司真正的机会，是第三阶段：产品内生数据，即产品运行过程中产生、过去不存在、外部无法直接购买的数据。

\begin{knowledgebox}{产品内生数据的例子}
\begin{tabularx}{\textwidth}{p{0.22\textwidth}p{0.34\textwidth}X}
\toprule
数据类型 & 怎么产生 & 为什么有壁垒 \\
\midrule
用户意图轨迹 & 用户在产品中不断描述、修正、选择 & 反映真实需求和偏好。 \\
生成失败样本 & 生成内容不满意时的修改过程 & 暴露模型短板和产品场景。 \\
Recipe & 可复用的生产方法和参数组合 & 不是单点内容，而是生产能力。 \\
商业闭环反馈 & 内容是否播放、付费、转化、复用 & 直接连接产品价值。 \\
\bottomrule
\end{tabularx}
\end{knowledgebox}

\subsection{本章小结}

数据三阶段解释了为什么应用公司不能只做薄 UI。真正的应用机会在于创造新数据：产品使用越多，越产生外部没有的意图、过程、失败、recipe 和商业反馈。

